% !PRINTED 394
\weberpage{394}
Denn man setze
\[
\begin{aligned}
N(\omega) &= a + 2^k b,\\
\omega_1 &= \omega (1 + 2^{k-1} x),
\end{aligned}
\]
worin $x$ rational. Dann ist $2k - 2 > k$, und folglich
\[
\begin{aligned}
N(\omega_1) &\equiv (a + 2^k b)(1 + 2^k x) \ (\mathrm{mod}\ 2^{k+1});\\
&\equiv a + 2^k (b + a x);
\end{aligned}
\]
wenn also
\[
a x + b \equiv 0 \ (\mathrm{mod}\ 2),
\]
so ist
\[
N(\omega_1) \equiv a \ (\mathrm{mod}\ 2^{k+1}).
\]
Aus alledem ergibt sich nun folgender Satz:

8. \gesperrt{Ist eine Zahl $r$ in $Z$ quadratischer Rest jeder in $D$ aufgehenden ungeraden Primzahl, und ist außerdem bei geraden $D$}
\[
\begin{aligned}
& \text{a)} \left\{
\begin{array}{ll}
r \equiv 1 \ (\mathrm{mod}\ 8) & \text{falls } \dfrac{D}{4} \equiv 0 \ (\mathrm{mod}\ 8),\\[4pt]
r \equiv +1, -1 \ (\mathrm{mod}\ 8) & \text{\glqq}\quad \dfrac{D}{4} \equiv 2 \ (\mathrm{mod}\ 8),\\[4pt]
r \equiv 1, 3 \ (\mathrm{mod}\ 8) & \text{\glqq}\quad \dfrac{D}{4} \equiv 6 \ (\mathrm{mod}\ 8),
\end{array}
\right. \\[6pt]
& \text{b)}\ r \equiv +1, -5 \ (\mathrm{mod}\ 4) \qquad \text{\glqq}\quad \dfrac{D}{4} \equiv 3, 4, 7 \ (\mathrm{mod}\ 8),
\end{aligned}
\]
\gesperrt{so ist $r$ Normenrest von jedem beliebigen zu $r$ teilerfremden Modul $m$.}

Diese Zahlen $r$ bilden eine Gruppe $R$, die wir die \gesperrt{Gruppe der absoluten Normenreste} nennen.

Ist $b$ ein fester quadratischer Nichtrest der ungeraden Primzahl $p$, und $z$ eine Zahl in $Z$, so ist entweder $z$ oder $b^{-1} z$ in der Gruppe der Normenreste enthalten. Ist also $P$ die Gruppe der Normenreste von $p^k$, so ist
\[
\begin{aligned}
(Z, P) &= 1, \quad \text{wenn } p \text{ nicht in } D \text{ aufgeht},\\
(Z, P) &= 2, \quad \text{wenn } p \text{ in } D \text{ aufgeht},
\end{aligned}
\]
ferner, wenn $L$ die Gruppe der Normenreste von $2^k$ ist:
\[
\begin{aligned}
(Z, L) &= 1, & D &\equiv 1 \ (\mathrm{mod}\ 4),\\
& & D &\equiv 4, 20 \ (\mathrm{mod}\ 32),\\
(Z, L) &= 2, & D &\equiv 8, 12, 16, 24, 28 \ (\mathrm{mod}\ 32),\\
(Z, L) &= 4, & D &\equiv 0 \ (\mathrm{mod}\ 32)\,^{1}).
\end{aligned}
\]

\par\noindent\rule{4cm}{0.4pt}\par
$^{1}$) Es ist nicht zu befürchten, daß das in \S\,99 und in Bd.\,II, \S\,21 erklärte Symbol $(Z, P)$ für den Index eines Teilers einer Gruppe mit dem in \S\,85 ähnlich bezeichneten erweiterten \gesperrt{Legendre-Jacobischen} Symbol verwechselt werde.
